% <<<PDFPAGE 31>>>
\chapter{推理服务系统架构}

\section{生产级推理服务器的组成}

生产环境下的大语言模型推理服务器是一个复杂的分布式系统，主要由以下组件构成。
\begin{itemize}
  \item \textbf{前端或 API 网关}（frontend / API gateway）。接收 HTTP 请求，对请求进行身份验证、限流，并将其加入队列。
  \item \textbf{调度器}（scheduler）。它是服务系统的大脑，决定当前批次运行哪些请求，同时管理显存和 SLO 约束。
  \item \textbf{执行引擎}（execution engine）。在 GPU 上实际执行前向传播并管理 KV 缓存块。
  \item \textbf{词元切分器}（tokenizer）。在文本与词元 ID 之间双向转换，通常在 CPU 上运行。
  \item \textbf{采样器}（sampler）。对 logits 应用温度、top-\(k\)、top-\(p\) 等采样规则。
\end{itemize}

\section{批处理策略}

\subsection{静态批处理}

最简单的做法是收集固定数量的请求，把它们的序列都填充至最长序列的长度，然后处理填充后的整批请求。这种静态批处理（static batching）的效率很低，序列长度差异越大，填充浪费的算力和显存就越多。

\subsection{连续批处理（迭代级批处理）}

Orca（Yu 等，2022）率先采用连续批处理，vLLM 和 TGI 都实现了这一策略。连续批处理以\emph{迭代}（解码步骤）为组批单位，突破了按请求固定组批的限制。每个解码步骤中，调度器都可以把已完成预填充的新请求加入正在运行的批次，同时立即移除已完成的序列。这消除了静态批处理的“长尾”问题，让 GPU 保持较高的利用率。

\subsection{分块预填充}

Sarathi-Serve（Agrawal 等，2024）发现，一个需要长时间预填充的请求可能独占 GPU 达数百毫秒，使解码阶段的请求得不到服务。这种现象称为“预填充抢占”（prefill piracy）。分块预填充（chunked prefill）将预填充计算分成较小的块，与解码步骤交错执行，从而使各类请求的时延更容易预测。

\section{预填充与解码分离}

预填充受计算能力限制，解码受访存带宽限制。两者需要的资源不同，可以分别交给针对各自工作负载优化的硬件。Splitwise 和 Mooncake 等分离式服务系统将两类节点物理分开。
% <<<PDFPAGE 32>>>
\begin{itemize}
  \item \textbf{预填充节点}。用大批次和长提示尽可能提高算力利用率。可以使用成本较低、针对计算优化的旧款 GPU。
  \item \textbf{解码节点}。优先获得尽可能大的访存带宽，支持高并发。采用每个模型参数可用 HBM 带宽较高的新款 GPU，通常能获益。
\end{itemize}

其中的难点是 \textbf{KV 缓存迁移}（KV cache migration）。预填充结束后，KV 缓存需要通过 NVLink、InfiniBand 或远程直接内存访问（remote direct memory access，RDMA），从预填充节点传到解码节点，增加了时延。月之暗面（Moonshot AI）旗下 Kimi 团队的 Mooncake 等系统在二级存储上构建分布式 KV 缓存池，将两类节点在缓存传输上的依赖完全解耦。

\section{考虑 SLO 的调度}

生产系统必须满足服务级别目标，通常设定 TTFT 和 TPOT 的 p95 或 p99 目标。朴素的先入先出（first in, first out，FIFO）调度可能使短请求排在长请求之后，长期得不到服务。考虑 SLO 的调度可以采用以下方法。
\begin{itemize}
  \item \textbf{优先级排队}（priority queuing）。让较短或优先级较高的请求排到前面。
  \item \textbf{抢占}。在长时间运行的请求仍有余量满足 SLO 时暂停它，释放其 KV 缓存块，转而服务优先级更高的请求。
  \item \textbf{请求路由}（request routing）。根据估计的请求长度和服务器负载，把请求送往不同的服务器池。
  \item \textbf{准入控制}（admission control）。系统过载时，拒绝请求或将其排队，以防无法满足 SLO。
\end{itemize}

\begin{interview}
\question{1}{请解释连续批处理。为什么它比静态批处理更适合大语言模型服务？}
静态批处理把长度相近的请求分在一组，对较短的请求做填充，然后处理该批次，直到所有请求都完成。如果一批有 32 个请求，其中一个的长度是其余请求的 10 倍，就会有 31 块 GPU 闲置，等待这一个请求结束。花在填充词元上的算力也无法收回。连续批处理（Orca、vLLM）以迭代为单位。在每个解码步骤，调度器检查哪些运行中的请求已经生成序列结束词元（end-of-sequence token，EOS），即 \texttt{<EOS>}，并立即用队列中新到的请求替换它们。长度不同的请求之间无须填充，各请求只占用实际所需的词元位置。这样无需全批等待最后完成的请求，GPU 利用率可一直保持在较高水平。按 Orca 论文的基准测试，时延相近时，连续批处理的吞吐量为静态批处理的 \(5\) 至 \(23\) 倍。

\translatornote{原书此处写作“31 块 GPU”。该例中的 32 是同批请求数。其余 31 个请求完成后空出的是批内处理位置，实际 GPU 数量取决于部署方式。}

\question{2}{什么是“预填充抢占”？分块预填充怎样处理这种情况？}
预填充抢占指一个带有很长提示的请求在预填充阶段独占 GPU，时间可能长达数百毫秒，导致批中所有处于解码阶段的请求停滞。预填充期间，GPU 对大词元矩阵执行 GEMM。
% <<<PDFPAGE 33>>>
此时正在生成词元的解码请求必须等待，使其 TPOT 骤增。分块预填充（Sarathi-Serve）把长提示的预填充划分为多个小块，例如每块 512 个词元，并将各块与解码步骤交错执行。这样，对 4096 个词元做预填充时，计算完 512 个词元就把算力让给解码批次。相较于整段预填充，解码的 TPOT 恶化幅度可降至原来的 \(1/10\) 至 \(1/5\)，同时只给总预填充时间带来很小的额外开销。

\question{3}{描述 vLLM 的架构。它有哪些主要创新与局限？}
vLLM（Virtual LLM）是高吞吐量服务系统，主要采用三项技术。

（1）\textbf{PagedAttention}。按非连续的块管理 KV 缓存，消除碎片。（2）\textbf{连续批处理}。按迭代组批，提高 GPU 利用率。（3）\textbf{定制注意力内核}。用 CUDA 内核执行分页注意力计算。

请求到达 API 服务器后，先做词元切分，再由 \texttt{LLMEngine} 调度。后者通过 \texttt{BlockManager} 管理 KV 块。定制执行引擎让请求在 GPU 工作进程上运行。

书中所述实现有以下局限。（1）调度器还不能完全按 SLO 调度，短请求可能长期得不到服务。（2）对于小模型，张量并行的通信开销可能很大。（3）大量使用 Python 的调度层会增加极短请求的时延。（4）基础版 vLLM 没有原生支持预填充与解码分离。（5）混合专家（Mixture of Experts，MoE）、非 Transformer 模型等非标准架构需要开发定制插件。

\translatornote{vLLM 0.8.0 已将预填充与解码分离列为实验性功能，0.8.5 的支持清单已收录 Mixtral 和 Mamba。见 \href{https://docs.vllm.ai/en/v0.8.0/features/disagg_prefill.html}{0.8.0 功能文档}和 \href{https://docs.vllm.ai/en/v0.8.5/models/supported_models.html}{0.8.5 模型清单}。}

\question{4}{分离式预填充与解码服务如何工作？实现起来有哪些工程难点？}
分离式服务将预填充任务路由至专用的预填充节点，将解码任务路由至专用的解码节点。请求到达路由器后，先发往预填充节点，由它执行提示的前向传播并产生 KV 缓存。缓存随后经 RDMA 或 InfiniBand 传给解码节点，后者生成词元直至请求结束。这样做面临几项工程挑战。（1）\textbf{KV 缓存传输时延}。通过 \(400\,\mathrm{Gb/s}\) 的 InfiniBand 链路传输 \(10\,\mathrm{GB}\) 的 KV 缓存约需 \(200\,\mathrm{ms}\)，往往比预填充本身还久。可逐层传输 KV 块，与解码形成流水线。（2）\textbf{资源平衡}。两类节点需要的算力与内存之比不同，任何一个节点池的资源配得过多或过少都会降低总体吞吐量。（3）\textbf{容错}。解码节点在生成中途发生故障时，必须恢复 KV 缓存。（4）\textbf{路由复杂性}。路由器必须估计请求长度，才能尽量均衡地分配负载。

\question{5}{你会用哪些指标来评价和比较两套大语言模型服务系统？}
我会比较以下指标。（1）\textbf{TTFT 的 p50、p95 和 p99}。衡量第一个词元出现得有多快，对交互式应用尤其重要。（2）\textbf{TPOT 的 p50 和 p95}。衡量后续词元到达的速度，反映流式输出是否顺畅。（3）\textbf{吞吐量（词元/秒）}。衡量给定硬件配置下系统的总输出速率。（4）\textbf{有效吞吐量}。同时满足 TTFT 和 TPOT 等所有 SLO 的请求所占的比例。（5）\textbf{GPU 利用率}。用模型浮点运算利用率（Model FLOP Utilization，MFU）来衡量，即峰值 GPU FLOP/s 中有多大比例实际用于模型计算。（6）\textbf{每百万词元的成本}。按输出量将总拥有成本（total cost of ownership，TCO）折算为每百万词元的成本。（7）\textbf{可扩展性}。衡量并发用户数增加时的性能变化。（8）\textbf{过载时的尾部时延}。观察系统是逐渐变慢，还是急剧恶化。我会选用贴近实际的请求分布来做基准测试，混合长短不一的提示和输出，避免只测人工构造的定长负载。
\end{interview}
